% -----------------------------------------------------------
%  Retrospectiva del Proyecto
% -----------------------------------------------------------

Este anexo presenta un análisis retrospectivo del proyecto DOMU, evaluando el grado de cumplimiento de los requerimientos, los riesgos que se materializaron y la comparación entre la planificación, el esfuerzo real y la estimación por puntos de casos de uso.

\section{Cumplimiento de Requerimientos por Módulo}

La Tabla~\ref{tab:cumplimiento} sintetiza el porcentaje de implementación alcanzado para cada módulo funcional del sistema, en relación con los requerimientos funcionales (RF) definidos en el Capítulo~\ref{cap:requerimientos}.

\begin{longtable}{|p{0.25\textwidth}|p{0.20\textwidth}|c|p{0.30\textwidth}|}
\caption{Cumplimiento de requerimientos por módulo}\label{tab:cumplimiento}\\
\hline
\textbf{Módulo} & \textbf{RFs cubiertos} & \textbf{\% Impl.} & \textbf{Observación} \\
\hline
\endfirsthead
\hline
\textbf{Módulo} & \textbf{RFs cubiertos} & \textbf{\% Impl.} & \textbf{Observación} \\
\hline
\endhead
\hline
\endfoot

Control de Accesos & RF\_01, RF\_03 & 100\% & Visitas, QR y pre-autorización operativos. \\
\hline
Encomiendas & RF\_04 & 100\% & Registro, seguimiento y notificación completos. \\
\hline
Tareas y Personal & RF\_05, RF\_06 & $\sim$70\% & Asignación y seguimiento implementados; CU\_08 (turnos) diferido a v2.0. \\
\hline
Finanzas & RF\_07, RF\_08 & $\sim$85\% & Gastos comunes, pagos y PDFs operativos; CU\_06 (conciliación bancaria) diferido. \\
\hline
Reservas & RF\_09 & 100\% & Amenidades, disponibilidad y reservas completos. \\
\hline
Votaciones & RF\_10 & 100\% & Creación, sufragio y cierre automático operativos. \\
\hline
Comunicación & RF\_11, RF\_17 & 100\% & Foro con categorías y chat WebSocket en tiempo real. \\
\hline
Tickets / Incidentes & RF\_12 & 100\% & Reporte, asignación y seguimiento de incidencias. \\
\hline
Proveedores & RF\_13 & 100\% & CRUD completo, órdenes y cotizaciones. \\
\hline
Mantenimiento & RF\_14 & 0\% & CU\_11 y CU\_18 diferidos íntegramente a v2.0. \\
\hline
Dashboards / KPI & RF\_15 & $\sim$50\% & Estadísticas de incidentes implementadas; dashboard financiero parcial. \\
\hline
RBAC & RF\_16 & 100\% & Cinco roles, permisos y aprobación de comunidades. \\
\hline
Marketplace & RF\_18 & 100\% & Publicación, edición e imágenes operativas. \\
\hline
Biblioteca & RF\_19 & 100\% & Documentos comunitarios con almacenamiento en Box. \\
\hline
Notificaciones RT & RF\_20 & 100\% & WebSocket de notificaciones con preferencias por usuario. \\

\end{longtable}

% === PLACEHOLDER: Reemplazar con imagen generada ===
% PROMPT para generador de imágenes (Gemini/GPT/DALL-E):
%
% "Genera un GRÁFICO DE BARRAS HORIZONTALES mostrando el porcentaje de implementación por módulo.
%
% NOTACIÓN:
% - GRÁFICO: barras horizontales, eje X = porcentaje (0% a 100%), eje Y = nombres de módulos.
% - BARRAS: rectangulares, sin bordes redondeados, relleno sólido.
% - COLORES: verde (#4CAF50) para 100%, naranja (#FF9800) para 50-99%, rojo (#F44336) para 0%.
% - ETIQUETAS: el porcentaje exacto al final de cada barra.
% - LÍNEA DE META: línea vertical punteada negra en 80% con etiqueta 'Meta mínima (80%)'.
%
% DATOS (15 módulos, de arriba a abajo):
% Control de Accesos: 100% (verde)
% Encomiendas: 100% (verde)
% Tareas/Personal: 70% (naranja)
% Finanzas: 85% (naranja)
% Reservas: 100% (verde)
% Votaciones: 100% (verde)
% Comunicación: 100% (verde)
% Tickets: 100% (verde)
% Proveedores: 100% (verde)
% Mantenimiento: 0% (rojo)
% Dashboards: 50% (naranja)
% RBAC: 100% (verde)
% Marketplace: 100% (verde)
% Biblioteca: 100% (verde)
% Notificaciones: 100% (verde)
%
% TÍTULO del gráfico: 'Porcentaje de implementación por módulo - DOMU v1.0'
% ESTILO: Fondo blanco, ejes en negro, sin grilla de fondo, tipografía sans-serif.
% Profesional para tesis universitaria. Etiquetas en español. Formato: PNG 300dpi."
\begin{figure}[H]
  \centering
  \fbox{\parbox{0.85\textwidth}{\centering\vspace{3cm}
  \textit{[Insertar gráfico: Porcentaje de implementación por módulo --- barras horizontales con codificación por color]}
  \vspace{3cm}}}
  \caption{Porcentaje de implementación por módulo del sistema DOMU.}
  \label{fig:cumplimiento-modulos}
\end{figure}

\section{Análisis de Éxito y Fracaso}

El proyecto DOMU alcanzó la implementación completa de \textbf{22 de los 27 casos de uso} definidos, lo que representa un \textbf{81\% de cumplimiento funcional}. Los cinco casos de uso diferidos a la versión~2.0 corresponden a funcionalidades de menor prioridad (conciliación bancaria, turnos de personal, mantenimiento preventivo, estacionamientos y bitácora de ascensores), cuya postergación fue una decisión deliberada para asegurar la calidad de los módulos entregados.

\textbf{Indicadores de éxito:}
\begin{itemize}
    \item Sistema funcional con más de 130 endpoints REST operativos y 2 canales WebSocket.
    \item 54 vistas frontend segmentadas por rol, cubriendo todos los flujos implementados.
    \item Arquitectura modular (Controller-Service-Repository) que facilita la extensión futura.
    \item Documentación de API automatizada mediante Swagger UI.
    \item Desviación de esfuerzo global de solo +3\% respecto a la planificación inicial.
\end{itemize}

\textbf{Aspectos a mejorar:}
\begin{itemize}
    \item La integración con Box API requiere renovación manual del token de desarrollador, limitando la autonomía del sistema en ambiente de producción.
    \item El módulo de dashboards quedó parcialmente implementado, ofreciendo solo estadísticas de incidentes.
    \item La pasarela de pago opera en modo simulado; la integración real con Mercado Pago queda pendiente para v2.0.
\end{itemize}

\section{Riesgos Materializados y sus Efectos}

De los ocho riesgos identificados en la etapa de planificación (ver Tabla~\ref{tab:riesgos} del Anexo~\ref{anexo:gestion}), dos se materializaron completamente y dos de forma parcial:

\begin{itemize}
    \item \textbf{R01 --- Complejidad del modelo financiero:} La necesidad de modelar 35~entidades en 9~dominios funcionales (frente a las $\sim$20 estimadas) extendió la fase de diseño de base de datos en aproximadamente una semana (+20\%). El impacto se absorbió parcialmente durante la implementación del frontend, que resultó más eficiente de lo previsto gracias a TanStack React Query.

    \item \textbf{R02 --- Postergación de funcionalidades:} La decisión de diferir 5~CUs se tomó en la semana~8 del desarrollo, tras evaluar que la implementación completa comprometería la calidad del software entregado. Esta decisión permitió dedicar más tiempo a la estabilización de los módulos financiero y de proveedores.

    \item \textbf{R05 --- Dependencia de Box API (parcial):} El token de desarrollador de Box tiene una vigencia limitada (60~minutos), lo que obligó a implementar un flujo de renovación manual durante el desarrollo. Para producción, se recomienda migrar a autenticación OAuth~2.0 con tokens de larga duración.

    \item \textbf{R07 --- Curva de aprendizaje (parcial):} La adopción de Javalin y TanStack React Query requirió un sprint inicial dedicado a prototipos. Sin embargo, la documentación oficial de ambas tecnologías resultó suficiente para alcanzar productividad plena a partir de la segunda semana.
\end{itemize}

\section{Comparación: Planificación, Esfuerzo Real y Estimación UCP}

La Tabla~\ref{tab:comparacion-retro} contrasta las tres fuentes de estimación del proyecto:

\begin{table}[H]
\centering
\begin{tabular}{|p{0.25\textwidth}|c|c|p{0.30\textwidth}|}
\hline
\textbf{Fuente} & \textbf{Horas} & \textbf{Desv. vs.\ Real} & \textbf{Método} \\
\hline
Carta Gantt & 340 & $-$3\% & Estimación por descomposición de fases \\
\hline
Esfuerzo real & 350 & --- & Medición directa por fase \\
\hline
UCP (LOE=20) & 5{,}218 & +1{,}391\% & Método estándar de Karner \\
\hline
UCP (LOE ajustado) & $\approx$350 & $\approx$0\% & LOE empírico = 1.34~hrs/UCP \\
\hline
\end{tabular}
\caption{Comparación de las tres fuentes de estimación.}
\label{tab:comparacion-retro}
\end{table}

\textbf{Análisis:}
\begin{itemize}
    \item La \textbf{carta Gantt} proporcionó la estimación más precisa (340~hrs vs.\ 350~hrs reales), con una desviación de apenas $-$3\%. Esto se debe a que la descomposición por fases se basó en la experiencia previa del equipo con tecnologías similares.

    \item El \textbf{método UCP estándar} (LOE=20~hrs/UCP) sobreestimó significativamente el esfuerzo, arrojando 5{,}218~horas. Esta divergencia es habitual en proyectos académicos donde el LOE calibrado para equipos empresariales no refleja las condiciones reales del desarrollo.

    \item Al aplicar un \textbf{LOE empírico de 1.34~hrs/UCP}, la estimación UCP converge con el esfuerzo real. Este valor es coherente con estudios que reportan LOE entre 1 y 5~hrs/UCP para proyectos académicos con alta reutilización de frameworks.

    \item Las \textbf{desviaciones internas} más significativas se dieron en diseño de BD (+20\%) y backend (+20\%), compensadas por frontend ($-$19\%) y PWA ($-$25\%), resultando en un balance global casi neutro.
\end{itemize}

\section{Conclusiones de la Retrospectiva}

\begin{enumerate}
    \item El proyecto \textbf{cumplió su objetivo general}: entregar un sistema funcional de administración de comunidades residenciales con cobertura del 81\% de los casos de uso planificados.

    \item Las \textbf{desviaciones de alcance} se manejaron mediante re-priorización oportuna (semana~8), comunicada a la profesora guía, sin afectar la fecha de cierre global.

    \item La \textbf{estimación por carta Gantt} resultó la más precisa para este tipo de proyecto, mientras que el método UCP requiere calibración del LOE para contextos académicos.

    \item La \textbf{arquitectura modular} adoptada facilitó la postergación de funcionalidades sin afectar los módulos implementados, validando la decisión de diseño.

    \item Para la \textbf{versión 2.0}, se recomienda priorizar: conciliación bancaria (CU\_06), turnos de personal (CU\_08) e integración real con pasarela de pagos, aprovechando la base arquitectónica ya establecida.
\end{enumerate}